\scanpage{102}

\fbox{Espaces de Schatten}

\textcircled{1} \uline{Normes de Schatten}

\uline{Définition}. \uline{Prop} Une norme de Schatten $N(x)$ $\unsure{\to|\text{est}}$ croissante pour $x$ positif. \unsure{\}|§|?}

\uline{Exemples} les $N_p$. Norme de Schatten définie par une Norme de S. sur $R^n$.

\uline{Définition} $N(x)$ pour $x\in\underline C^{\underline N}$. $\ell^N\ni x\Longleftrightarrow N(x)$ fini, $N$ est une norme sur $\ell^N$. Espace $\ell_N$, adhérence de $C^{(\underline N)}$ dans $\ell^N$

\uline{Inclusions} pour $N'\leq N''$, en particulier

\[
\ell^1\subset\ell^N\subset\ell^\infty,\qquad \ell^1\subset\ell_N\subset c_0
\]

1 \uline{Prop} $\ell_N\subset\underline c_0$ sauf si $N$ est équivalente à $N_\infty$

\uline{Polarité}, $(N^\circ)^\circ=N^\circ$\,: \uline{Prop} le dual de $\ell_N$ est $\ell^{N^\circ}$. \uline{Corollaire} $\ell^N$ est complet, $\ell_N$ est complet. \uline{Prop} $\ell_N$ réflexif $\Longleftrightarrow\ell^N$ réflexif $\Longleftrightarrow\ell_{N^\circ}$ réflexif $\Longleftrightarrow\ell^{N^\circ}$ réflexif $\Longleftrightarrow\ell_N=\ell^N$ et $\ell_{N^\circ}=\ell^{N^\circ}$.

\deleted{\unsure{Exemples}} \uline{Cas des $N_p$} On a $(N_p)^\circ=N_{\!p'}^{\deleted{\unsure{p}}}$. On retrouve la réflexivité des $\ell^p$ $(1<p<+\infty\unsure{)}$

\textcircled{2} \uline{Espaces de Schatten}

\deleted{\uline{Définition} Si $u\in L_0^{\unsure{N}}(E,F)$, \inserted{$N(u)=$}$N((\rho_i(u)))$\unsure{;} $u\in L^N(E,F)\Longleftrightarrow N(u)<+\infty$, \unsure{Prop} $N(u)$ est une norme sur $L^N(E,F)$, $L_N(E,F)$ adhérence de $E'\otimes F$ ($N$ non équivalente à $N_\infty$). $N(u)=N(u^*)$, $N(uv)\leq\lVert v\rVert N(u)$, $N(vu)\leq\lVert v\rVert N(u)$.}

\deleted{\uline{Définition} (Si $N$ équivalente à $N_\infty$) $N(u)=\underset{\substack{N^\circ(v)\leq 1\\v\in F'\otimes E}}{\operatorname{Sup}}|\operatorname{Tr}vu|$, Si $u\in L_0(E,F)$, on retrouve $N(u)=N((\rho_i(u)))$, i.e, pour tout $N$}

\[
\deleted{\unsure{(}N((\rho_i(u)))=\underset{\substack{v\in F'\otimes E\\N^\circ(v)\leq 1}}{\operatorname{Sup}}|\operatorname{Tr}vu|}
\]

\uline{Définition\,:} Si $u\in L_0(E,F)$, posons $N(u)=N(\rho_i(u))$

\[
N(u)=\deleted{N((\rho_i(u)))}=\underset{\substack{v\in F'\otimes E\\\deleted{\unsure{N^\circ((\rho_i(v)))\leq 1|N^\circ(v)\leq1}}\\N^\circ(v)}}{\operatorname{Sup}}|\operatorname{Tr}vu|
\]

Définissons, pour $u\in L(E,F)$, $N(u)$ par la formule ci-dessus. Soit $L^N(E,F)\ni u\Longleftrightarrow\deleted{\unsure{N(u)=\operatorname{Sup}\operatorname{Tr}vu}}N(u)<+\infty$, $N\inserted{(u)}$ \unsure{est} une \uline{norme} sur $L^N(E,F)$ \deleted{et on a} soit $L_N(E,F)$ l'adhérence de $E'\otimes F$ dans $L^N(E,F)$

\deleted{identique aux} \uline{Prop} $u\in L_N(E,F)\Longleftrightarrow u\in L_0(E,F)$ et $(\rho_i(u))\in\ell_N$ \deleted{\unsure{Cor}} \deleted{\uline{Prop}} \uline{Inclusions} pour $N'\leq N''\Rightarrow\deleted{\unsure{\ell^{N''}\subset\ell^{N'}}}\,L^{N''}(E,F)\subset L^{N'}(E,F)$ et $L_{N''}(E,F)\subset L_{N'}(E,F)$, en particulier

\[
\begin{aligned}
L^1(E,F)&\subset L^N(E,F)\subset L(E,F)\\
L^1(E,F)&\subset L_N(E,F)\subset L_0(E,F)
\end{aligned}
\]

\deleted{\uline{Prop} Si $N$ non équivalente à $N_\infty$, alors}

\[
\text{Prop}\quad\boxed{N(u)=N(u^*)\quad N(uv)\leq\lVert v\rVert N(u),\ N(vu)\leq\lVert v\rVert N(u)}
\]

\deleted{\uline{Caractérisation} \unsure{$N(u)$, utilisant.}} (par transposition il suffit de le voir pour $u\in L_0(E,F)$, celà résulte de $\rho_i(vu)\leq\lVert v\rVert\rho_i(u)$, \ldots)
% archive sheet number: 103
